% ============================================================
% Sección 4 — Ejecución del caso de la profesora
% Todas las cifras salen de results/caso_entrada_G2_red_15_nodos/ (T06, semilla 20261105).
% ============================================================
\section{Ejecución del caso de la profesora}\label{sec:caso}

Esta sección reporta lo que produjo el código sobre \texttt{project-docs/entrada\_G2\_red\_15\_nodos.json} con
\texttt{python -m ada\_g2 run}. \textbf{Nada de lo que sigue se calculó a mano}: cada cifra se toma de uno de los cuatro archivos de
\texttt{results/caso\_entrada\_G2\_red\_15\_nodos/} (\texttt{resultado.json}, \texttt{tiempo.csv}, \texttt{traza.md}, \texttt{traza.csv}) y la fuente se indica
junto a cada tabla. Salvo que se diga otra cosa, las cifras reales se muestran con cuatro decimales. Parámetros de la corrida
(\texttt{resultado.json}, campo \texttt{parametros}): $N=15$, $T=3$, $k=5$, $B=30$, $L=24$, $105$ aristas, modo \texttt{ambos}
(cupo $k$ y costo $L$ a la vez, D-03), $q=0$ (D-07), política de gasto \texttt{concentrado} (D-22), semilla $20261105$. Entorno
(\texttt{resultado.json}, campo \texttt{entorno}): Python 3.12.10, Windows-11-10.0.26300-SP0. El campo \texttt{avisos} está vacío y
\texttt{todo\_cumple} vale \texttt{true}.

% ------------------------------------------------------------
\subsection{Resultado por ronda}

\begin{table}[h]
\centering
\small
\begin{tabular}{@{}c l c c l c c@{}}
\toprule
Ronda & Selección $S_t$ & $f(S_t)$ & $\sum c_i$ & Reparto (N09/N11/N13/N14/N15) & Topes efectivos & $\sum_i u_i(x_i)$\\
\midrule
1 & N09, N11, N13, N14, N15 & 145.5795 & 23 & 7 / 5 / 5 / 7 / 6 & 18 / 18 / 18 / 18 / 18 & 88.5640\\
2 & N09, N11, N13, N14, N15 & 145.5795 & 23 & 7 / 5 / 5 / 7 / 6 & 11 / 13 / 13 / 11 / 12 & 88.5640\\
3 & N09, N11, N13, N14, N15 & 145.5795 & 23 & 4 / 8 / 8 / 4 / 6 & 4 / 8 / 8 / 4 / 6 & 85.6108\\
\midrule
\multicolumn{2}{l}{Totales} & $\sum_t f=436.7386$ & & & & $\sum_t u=262.7387$\\
\bottomrule
\end{tabular}
\caption{Resultado por ronda. Origen: \texttt{resultado.json}, campos \texttt{rondas[].seleccion}, \texttt{f}, \texttt{costo\_participacion}, \texttt{reparto}, \texttt{topes\_efectivos}, \texttt{utilidad} y \texttt{totales.suma\_f}, \texttt{totales.suma\_utilidad}.}\label{tab:rondas}
\end{table}

\paragraph{Lectura.} (i) Con $q=0$ y ningún nodo agotado, el criterio de selección no cambia entre rondas (D-21): $S_1=S_2=S_3$. Solo cambia el reparto, porque los topes restantes $r_i(t)$ bajan (D-06). (ii) Los topes \emph{sí} son activos en la ronda 3: tras las rondas 1 y 2 cada participante ha gastado $7+7=14$, $5+5=10$, $5+5=10$, $7+7=14$ y $6+6=12$ unidades, le quedan $4,8,8,4,6$ (suman $30=B$) y el reparto de la ronda 3 coincide con esos topes. En esa ronda no hay nada que optimizar sobre la repartición: toda la capacidad restante se usa. (iii) Cada uno de los cinco participantes termina con gasto total $18$, igual a su \texttt{presupuesto\_total\_nodo} (\texttt{resultado.json}, \texttt{totales.gasto\_por\_nodo} y \texttt{totales.presupuesto\_total\_nodo}).

% ------------------------------------------------------------
\subsection{Traza resumida de la ronda 1}

\paragraph{Selección (Parte 2), heurística elegida: \texttt{voraz\_ganancia}.} Origen: \texttt{traza.md}, sección «Ronda 1 / Selección», con seis decimales como allí. La última columna da el mejor competidor en ese paso, es decir, el segundo mayor de las ganancias marginales de los candidatos factibles (columna «ganancias de los candidatos factibles» de \texttt{traza.md}).

\begin{table}[h]
\centering
\small
\begin{tabular}{@{}c l r c r l@{}}
\toprule
Paso & Elegido & Ganancia marginal $\Delta(i\mid S)$ & Costo acum. & $f$ acum. & Mejor competidor\\
\midrule
1 & N11 & 103.490000 & 4  & 103.490000 & N09 (95.7000)\\
2 & N14 & 24.278250  & 7  & 127.768250 & N09 (23.3850)\\
3 & N09 & 9.932737   & 10 & 137.700987 & N07 (9.3122)\\
4 & N15 & 5.530613   & 16 & 143.231600 & N13 (4.1325)\\
5 & N13 & 2.347920   & 23 & 145.579520 & N12 (2.2442)\\
\bottomrule
\end{tabular}
\caption{Las cinco elecciones del voraz en la ronda 1. Las ganancias decrecen paso a paso (submodularidad, Lema~\ref{lem:submodular}) y la suma de las ganancias marginales es $f(S)$.}\label{tab:traza-sel}
\end{table}

Las otras dos heurísticas del modo \texttt{ambos} dieron (\texttt{resultado.json}, \texttt{rondas[0].detalle\_seleccion.alternativas}): \texttt{voraz\_razon}, $S=\{$N05, N09, N14, N11, N15$\}$ con $f=145.1056$ y costo $18$; y \texttt{singleton}, $S=\{$N11$\}$ con $f=103.4900$ y costo $4$. Se elige la de mayor $f$ (empate: la primera en el orden ganancia, razón, singleton; D-23). El voraz por ganancia termina con $k=5$ nodos elegidos (cupo agotado) y costo $23\le24$.

\paragraph{Reparto (Parte 1), unidad a unidad.} Participantes N09, N11, N13, N14, N15; tope de cada uno $18$. Origen: \texttt{traza.md}, sección «Reparto (Parte 1), unidad a unidad».

\begin{table}[h]
\centering
\small
\begin{tabular}{@{}c l c r r@{}}
\toprule
Unidad & Nodo & $x$ antes & Ganancia marginal $g_i(x)$ & Utilidad acum.\\
\midrule
1 & N14 & 0 & 7.000787 & 7.000787\\
2 & N09 & 0 & 6.862157 & 13.862944\\
3 & N15 & 0 & 6.584898 & 20.447842\\
4 & N13 & 0 & 5.753122 & 26.200963\\
5 & N11 & 0 & 5.267919 & 31.468882\\
$\vdots$ & & & &\\
30 & N09 & 6 & 1.321961 & 88.563953\\
\bottomrule
\end{tabular}
\caption{Primeras cinco unidades y última unidad del reparto de la ronda 1. La última ganancia, $1.3220$, es el multiplicador $\lambda$ del certificado de dualidad (Sección~\ref{sec:certificados}).}\label{tab:traza-rep}
\end{table}

El reparto final es N09$=7$, N11$=5$, N13$=5$, N14$=7$, N15$=6$ con utilidad $88.563953$. El montículo saca siempre la mayor ganancia marginal disponible; como $g_i(x)=a_i\ln\frac{x+2}{x+1}$ es decreciente en $x$, las 30 unidades tomadas son las 30 mayores ganancias del multiconjunto (Teorema~\ref{teo:optimo}).

% ------------------------------------------------------------
\subsection{Verificación de restricciones}\label{sec:verificacion}

El código comprueba explícitamente, por ronda, las siete restricciones y la coherencia numérica de la Tabla~\ref{tab:verif}, y al final la restricción global del presupuesto total por nodo. Origen: \texttt{resultado.json}, campo \texttt{rondas[].verificacion\_restricciones} y \texttt{verificacion\_global}; la columna «valor» corresponde a la ronda 1.

\begin{table}[h]
\centering
\small
\begin{tabular}{@{}l l r r c@{}}
\toprule
Restricción & Significado & Valor & Límite & Cumple\\
\midrule
$|S|\le k$ & cupo de participantes & 5 & 5 & sí\\
$\sum c\le L$ & costo de participación & 23 & 24 & sí\\
$\sum x\le B$ & presupuesto de ruido de la ronda & 30 & 30 & sí\\
$x_i\le$ tope$_i$ & tope por nodo (política) & $-11$ & 0 & sí\\
$x_i\le$ restante$_i$ & presupuesto total restante & (sin violaciones) & (ninguna) & sí\\
$x_i\ge0$ entero & dominio & 5 (mínimo de $x$) & 0 & sí\\
$S\subseteq$ disponibles & participantes disponibles & (sin violaciones) & (ninguna) & sí\\
$f$ incremental $=f$ directa & coherencia numérica & $2.8422\times10^{-14}$ & $10^{-9}$ & sí\\
\bottomrule
\end{tabular}
\caption{Restricciones verificadas en la ronda 1. En las rondas 2 y 3 todas también cumplen; en la ronda 3 el valor de «tope por nodo» es $0$ (holgura nula: el tope es activo) y el mínimo de $x$ es $4$. Restricción global $\sum_t x_{i,t}\le R_i$: sin violaciones; \texttt{todo\_cumple}$=$\texttt{true}.}\label{tab:verif}
\end{table}

En la fila de tope por nodo, «valor» es el máximo de $x_i-\text{tope}_i$ entre los participantes: $-11$ significa que el nodo más cercano a su tope todavía tiene $11$ unidades de holgura.

% ------------------------------------------------------------
\subsection{Certificados y comparación con el exacto}\label{sec:certificados}

\paragraph{Certificados por ejecución (D-14; demostrados en el código, no asumidos).} Origen: \texttt{resultado.json}, \texttt{certificado\_reparto} y \texttt{certificado\_seleccion}.
\begin{itemize}[leftmargin=*]
\item \emph{Parte 1, cota de dualidad de Lagrange.} Rondas 1 y 2: $\lambda=1.3220$; ronda 3: $\lambda=0.8952$. En las tres rondas la cota superior coincide con la utilidad obtenida ($88.5640$, $88.5640$ y $85.6108$) y la \textbf{brecha de dualidad es $0$} (tolerancia $10^{-9}$): el reparto es óptimo sin necesidad de conocer el óptimo por otro método.
\item \emph{Parte 2, cota a posteriori por programación lineal.} En las tres rondas $f(S)=145.5795$, el programa lineal acota lo que podría añadir cualquier otro conjunto factible en $6.4076$, y la \textbf{cota superior del óptimo es $151.9871$}. La \textbf{razón certificada es $0.9578$}: $f(S)\ge0.9578\cdot\OPT$, con las restricciones de cupo y costo a la vez. Este certificado vale aunque el modo \texttt{ambos} no tenga garantía a priori (R-01).
\end{itemize}

\paragraph{Comparación con el exacto (referencia de validación, no parte del algoritmo).} Origen: \texttt{resultado.json}, \texttt{rondas[].comparacion\_exacto}.
\begin{itemize}[leftmargin=*]
\item \emph{Modo \texttt{ambos} (el de la ejecución).} La fuerza bruta evalúa $3\,893$ subconjuntos factibles de los $4\,944$ posibles de tamaño $\le5$ (contando el vacío; $4\,943$ no vacíos, como en la Sección~\ref{sec:supuestos}). El óptimo es $S^\ast=\{$N09, N11, N13, N14, N15$\}$ con $f^\ast=145.5795$, el mismo conjunto del voraz: \textbf{razón voraz/óptimo $=1.0000$}. En las tres rondas el reparto del voraz coincide con el de la programación dinámica: $88.5640$, $88.5640$ y $85.6108$ (diferencia absoluta $1.4211\times10^{-14}$, \texttt{reparto\_optimo}$=$\texttt{true}).
\item \emph{Modo \texttt{k} (solo cupo, sin costo).} Cifras de \texttt{tests/test\_caso\_profesora.py} (reproducibles con \texttt{--modo k}; no forman parte de los archivos de \texttt{results/} porque la corrida entregada es en modo \texttt{ambos}): el óptimo con $|S|\le5$ es $S^\ast=\{$N09, N11, N12, N13, N15$\}$ con $f^\ast=145.9072$, explorando $4\,944$ subconjuntos, y el voraz ($f=145.5795$) logra la razón $0.99775$ ($=145.5795/145.9072$). La garantía teórica en este modo es $1-(1-1/5)^5=2101/3125=0.67232\ge1-1/e\approx0.6321$ (Teorema~\ref{teo:nwf}); la razón observada la supera con holgura.
\end{itemize}
Una razón de $1.0000$ en un único caso \textbf{no es una garantía}: es una medición. La garantía es la del Teorema~\ref{teo:nwf}, solo en el modo \texttt{k}.

% ------------------------------------------------------------
\subsection{Tiempos de ejecución}\label{sec:tiempos}

Medición con \texttt{time.perf\_counter}, un calentamiento sin medir y $5$ repeticiones (D-11); se reporta la mediana. Origen: \texttt{tiempo.csv} (columna \texttt{mediana\_s}, convertida a milisegundos sin redondeo adicional).

\begin{table}[h]
\centering
\small
\begin{tabular}{@{}l r@{\hspace{2em}}l r@{}}
\toprule
Fase & Mediana (ms) & Fase & Mediana (ms)\\
\midrule
\texttt{carga\_json} & 0.5687 & \texttt{seleccion\_r1} & 0.4073\\
\texttt{seleccion} (3 rondas) & 1.2734 & \texttt{reparto\_r1} & 0.0433\\
\texttt{reparto} (3 rondas) & 0.1350 & \texttt{seleccion\_r2} & 0.4279\\
\texttt{verificacion} & 0.0526 & \texttt{reparto\_r2} & 0.0416\\
\texttt{certificados} & 5.2023 & \texttt{seleccion\_r3} & 0.3949\\
\texttt{pipeline\_total} & 6.6853 & \texttt{reparto\_r3} & 0.0484\\
\midrule
\multicolumn{4}{l}{\texttt{referencia\_exacta} (1 sola corrida, fuerza bruta + DP de las 3 rondas): 51.2363 ms}\\
\bottomrule
\end{tabular}
\caption{Tiempos medianos del caso de la profesora. Los certificados (programa lineal incluido) dominan el pipeline; el algoritmo de selección y de reparto, sin certificados, cuesta $1.2734+0.1350=1.4084$ ms.}\label{tab:tiempos}
\end{table}

Con una sola instancia ($N=15$) estos tiempos no permiten estimar la curva $T(n)$; esa estimación requiere el barrido por tamaños de la validación experimental (T13), que aún no se ha ejecutado. Se incluyen aquí solo como evidencia de que el caso se resuelve en pocos milisegundos.

% ------------------------------------------------------------
\subsection{Advertencia de saturación (R-02)}\label{sec:saturacion}

La instancia de la profesora es un grafo completo con pesos $\ge0.2$. Como $f(S)=\sum_j a_j[1-\prod_{i\in S}(1-w_{ij})]$ y cada nodo queda cubierto con probabilidad cercana a $1$ por cualquier conjunto de cinco, $f$ \textbf{se satura}: $f(V)=\sum_ja_j=147.8$ (suma de \texttt{coeficiente\_utilidad} en \texttt{project-docs/entrada\_G2\_red\_15\_nodos.json}) y la selección de la ronda 1 alcanza $145.5795/147.8=0.9850$ de ese máximo. Consecuencias que conviene decir con claridad:
\begin{itemize}[leftmargin=*]
\item Las tres heurísticas del modo \texttt{ambos} dan valores cercanos (\texttt{voraz\_razon}: $145.1056$, diferencia $0.4740$ frente al voraz por ganancia), y el óptimo coincide con el voraz. El caso \textbf{discrimina poco} entre selecciones.
\item Por eso este caso no es evidencia fuerte de la calidad de la aproximación. La verificación de la cota $1-1/e$ se hace sobre familias sintéticas (Sección~\ref{sec:cota}, párrafo de verificación empírica) y sobre la instancia ajustada de D-15, no sobre este grafo.
\end{itemize}
